\unnumberedSection{This Book}

This book will describe my version of a layout control system with hardware and software designed from the ground up. The big question is why build one yourself. Why yet another one? There is after all no shortage on such systems readily available. And there are great communities out there already underway. The key reason for doing it yourself is that it is simply fun and you learn a lot about standards, electronics and programming by building a system that you truly understanding from the ground up. To say it with the words of Richard Feynman

\begin{quotation}
    \textit{ "What I cannot create, I do not understand. -- Richard Feynman"}
\end{quotation}

Although it takes certainly longer to build such a system from the ground up, you still get to play with the railroad eventually. And even after years, you will have a lay  out control system properly documented and easy to support and enhance further. Not convinced? Well, at least this book should be interesting and give some ideas and references how to go after building such a system.

\unnumberedSection{Parts and Chapters}

// ??? rework...

The book is organized into several parts and chapters. The first part describes the underlying concepts of the layout control system. Hardware modules, nodes, ports and events and their interaction are outlined. Next, the set of message that are transmitted between the components and the message protocol flow illustrate how the whole system interacts. With the concepts in place, the software library available to the node firmware programmer is explained along with example code snippets. After this section, we all have a good idea how the system configuration and operation works. The section is rounded up with a set of concrete programming examples.

Perhaps the most important part of a layout control system is the management of locomotives and track power. After all, we want to run engines and play. Our system is using the DCC standard for running locomotives and consequently DCC signals need to be generated for configuring and operating an engine. A base station module will manage the locomotive sessions, generating the respective DCC packets to transmit to the track. Layouts may consist of a number of track sections for which a hardware module is needed to manage the track power and monitor the power consumption. Finally, decoders can communicate back and track power modules need to be able to detect this communication. Two chapters will describe these two parts in great detail.

The next big part of the book starts with the hardware design of modules. First the overall outline of a hardware module and our approach to module design is discussed. Building a hardware module will rest on common building blocks such as a CAN bus interface, a microcontroller core, H-Bridges for DCC track signal generation and so on. Using a modular approach the section will describe the building blocks developed so far. It is the idea to combine them for the purpose of the hardware module.

With the concepts, the messages and protocol, the software library and the hardware building blocks in place, we are ready to actually build the necessary hardware modules. The most important module is the base station. Next are boosters, block controllers, handhelds, sensor and actor modules, and so on. Finally, there are also utility components such as monitoring the DCC packets on the track, that are described in the later chapters. Each major module is devoted a chapter that describes the hardware building blocks used, additional hardware perhaps needed, and the firmware developed on top of the core library specifically for the module. Finally, there are several appendices with reference information and further links and other information.

\unnumberedSection{A final note}

A final note. "Truly from the ground up" does not mean to really build it all yourself. As said, there are standards to follow and not every piece of hardware needs to be built from individual parts. There are many DCC decoders available for locomotives, let's not overdo it and just use them. There are also quite powerful controller boards along with great software libraries for the micro controllers, such as the CAN bus library for the AtMega Controller family, already available. There is no need to dive into all these details.

The design allows for building your own hardware just using of the shelf electronic components or start a little more integrated by using a controller board and other breakout boards. The book will however describe modules from the ground up and not use controller boards or shields. This way the principles are easier to see. The appendix section provides further information and links on how to build a system with some of the shelf parts instead of building it all yourself. With the concepts and software explained, it should not be a big issue to build your own mix of hardware and software.

I have added most of the source files in the appendix for direct reference. They can also be found also on GitHub. ( Note: still to do... ) Every building block schematic shown was used and tested in one component or another. However, sometimes the book may not exactly match the material found on the web or be slightly different until the next revision is completed. Still, looking at portions of the source in the text explain quite well what it will do. As said, it is the documentation that hopefully in a couple of years from now still tells you what was done so you can adapt and build upon it. And troubleshoot.

The book hopefully also helps anybody new to the whole subject with good background and starting pointers to build such a system. I also have looked at other peoples great work, which helped a lot. What I however also found is that often there are rather few comments or explanations in the source and you have to partially reverse engineer what was actually build for understanding how things work. For those who simply want to use an end product, just fine. There is nothing wrong with this approach. For those who want to truly understand, it offers nevertheless little help. I hope to close some of these gaps with a well documented layout system and its inner workings.

In the end, as with any hobby, the journey is the goal. The reward in this undertaking is to learn about the digital control of model railroads from running a simple engine to a highly automated layout with one set of software and easy to build and use hardware components. Furthermore, it is to learn about how to build a track signaling system that manages analog and digital engines at the same time. So, enjoy.



